% cap1.tex

\chapter{Introducción}
\label{ch:intro} % la etiqueta para referencias

{\color{blue} En esta sección de debe incluir el fundamento general de la motivación de la investigación desarrollada en la tesis. También se debe incluir la Hipótesis de trabajo, Objetivo General y específicos. (Máximo 5 páginas). Usar mismo tamaño, tipo de letra y espacio del este formato).}

\vspace{0.5cm}

The \Gls{latex} typesetting markup language is specially suitable 
for documents that include \gls{maths}. \Glspl{formula} are 
rendered properly an easily once one gets used to the commands.

Given a set of numbers, there are elementary methods to compute 
its \acrlong{gcd}, which is abbreviated \acrshort{gcd}. This 
process is similar to that used for the \acrfull{lcm}.


\section{Hipótesis}


\section{Objetivos}

Objetivo General

Objetivos específicos


\section{Estructura del informe}

En este capítulo \ref{ch:intro} revisamos.... En el próximo Capítulo \ref{ch:litRev} veremos... En el Capítulo \ref{ch:metodo} ...

{\color{blue} (2 a 3 líneas que informe al lector el contenido de cada capítulo) }